\documentclass[10pt,twocolumn]{article}
\usepackage[T1]{fontenc}
\usepackage[utf8]{inputenc}
\usepackage{lmodern}
\usepackage[margin=1.8cm,top=2.4cm,bottom=2cm,columnsep=0.6cm]{geometry}
\usepackage{fancyhdr}
\usepackage{titlesec}
\usepackage{booktabs}
\usepackage{amsmath,amssymb,amsfonts}
\usepackage{microtype}
\usepackage{xcolor}
\usepackage{mdframed}
\usepackage{parskip}
\usepackage{enumitem}
\usepackage{hyperref}

\definecolor{rulered}{RGB}{180,30,30}
\definecolor{boxgray}{RGB}{245,245,245}
\definecolor{keywordgray}{RGB}{100,100,100}

\pagestyle{fancy}
\fancyhf{}
\fancyhead[L]{\textsc{\small Claude Mag}}
\fancyhead[C]{\small\textit{Machine Learning Research Digest}}
\fancyhead[R]{\small 2026-08-05}
\fancyfoot[C]{\small\thepage}
\renewcommand{\headrulewidth}{0.8pt}

\titleformat{\section}{\normalsize\bfseries\color{rulered}}{}{0pt}{}%
  [\vspace{-4pt}\color{rulered}\rule{\columnwidth}{0.4pt}]
\titlespacing{\section}{0pt}{8pt}{4pt}

\hypersetup{colorlinks=true,urlcolor=rulered,linkcolor=black}

\begin{document}
\setlength{\parindent}{0pt}
\setlength{\parskip}{4pt}

{\small\color{keywordgray}\textbf{Keywords:}
  discrete diffusion \textbullet{}
  mixture of experts \textbullet{}
  fast text generation \textbullet{}
  open-weight LLM}\par\vspace{4pt}

{\LARGE\bfseries\raggedright
  DiffusionGemma Technical Report}\par\vspace{4pt}

{\large\itshape\raggedright
  Google DeepMind's Open 26B MoE Model Generates Text at 1000+ Tokens per Second
  via Block-Parallel Discrete Diffusion}\par\vspace{6pt}

{\small\textbf{By} Adrien Ali Ta\"iga et al.\ (DiffusionGemma Team,
  Google DeepMind)}\par\vspace{2pt}

{\small\color{keywordgray}arXiv:2608.00146 \textbullet{} August 2026}
\par\vspace{2pt}\noindent{\color{rulered}\rule{\columnwidth}{0.6pt}}\vspace{6pt}

Autoregressive text generation has a fundamental throughput ceiling: it emits one token
per forward pass, making latency proportional to sequence length. \textit{DiffusionGemma}
eliminates this bottleneck by iteratively denoising an entire 256-token canvas in parallel,
achieving over 1000 tokens per second on a single NVIDIA H100---up to $4\times$ faster than
comparable AR models---while fitting within 18\,GB VRAM when quantised and released
under the Apache 2.0 licence.

\begin{mdframed}[backgroundcolor=boxgray,linecolor=rulered,linewidth=1pt,
    innerleftmargin=6pt,innerrightmargin=6pt,
    innertopmargin=6pt,innerbottommargin=6pt]
\textbf{\small AT A GLANCE}\par\vspace{4pt}
\begin{description}[leftmargin=0pt,labelsep=4pt,itemsep=2pt,parsep=0pt,
    font=\small\bfseries]
  \item[Problem:] Autoregressive LLMs generate one token per forward pass, making
    throughput fundamentally linear in sequence length and bottlenecked by memory bandwidth.
  \item[Why it matters:] A $4\times$ throughput gain at matched quality halves the
    inference cost of frontier LLMs, directly enabling larger, cheaper deployments.
  \item[Method:] Fine-tune the Gemma 4 26B MoE backbone with (1) SFT for bidirectional
    denoising and (2) RL + sampler distillation for joint quality and efficiency gains.
  \item[Result:] 1000+ tok/s on H100, 700+ tok/s on RTX~5090; fits in 18\,GB VRAM when
    quantised; open-weights under Apache 2.0.
  \item[Evidence:] Speed benchmarks on NVIDIA H100 and RTX~5090; quality evaluated against
    equivalent-parameter AR models.
\end{description}
\end{mdframed}

\section{The Big Picture}

Large language models generate text by sampling one token at a time in a strict
left-to-right order. This autoregressive (AR) constraint is simple and ensures
the model can condition each new token on all prior context, but it imposes a hard
throughput ceiling: no matter how wide the model is, each forward pass produces
exactly one token. For a sequence of $N$ tokens, the model must run $N$ serial
passes, making decode time $O(N)$ in the number of forward passes regardless of
hardware parallelism.

Diffusion-based text generation breaks this constraint by treating generation as
iterative denoising over a complete sequence. DiffusionGemma represents the state
of the art on this front: a 26B-parameter open-weights model that generates full
blocks of tokens in parallel, with measured throughput exceeding 1000 tokens per
second on a single H100---a $4\times$ improvement over equivalent AR models.

\section{Why Existing Approaches Fall Short}

Earlier attempts at diffusion-based language models (e.g.\ MDLM, SEDD) trained
from scratch on text corpora and did not match the quality of AR models of similar
parameter count. The quality gap arose because diffusion LMs require learning
bidirectional context modelling, a fundamentally harder pre-training problem than
the unidirectional next-token prediction that gives AR models their strong priors.

Prior fast-generation methods for AR models (speculative decoding, token streaming)
provide modest speedups (2--3$\times$) but remain bounded by the sequential
auto-regressive paradigm. DiffusionGemma bypasses this bound entirely.

\section{Core Mechanism}

\textbf{Architecture.}
DiffusionGemma is based on Gemma~4, specifically the 26B-A4B Mixture-of-Experts
variant: 26B total parameters with only 3.8B activated per token during inference.
The sparse MoE routing means most expert parameters are dormant in any given forward
pass, keeping latency low while retaining the model's large capacity.

\textbf{Block-Autoregressive Multi-Canvas Sampling.}
At inference, DiffusionGemma begins with a canvas of 256 \texttt{[MASK]} tokens.
Each denoising iteration:
\begin{enumerate}[itemsep=2pt,parsep=0pt]
  \item Predicts the probability distribution over the full vocabulary for every
    masked position simultaneously;
  \item Samples and commits the most confident subset of positions;
  \item Repeats until the canvas is fully populated.
\end{enumerate}
Once committed, positions are not revised. In practice far fewer than 256 iterations
are required, as many positions are filled confidently in the first few steps.

\section{Technical Formulation}

Let $\mathbf{x} = (x_1, \ldots, x_L)$ be a sequence of $L$ tokens.
At each denoising step $t$ the model receives a partially masked sequence
$\tilde{\mathbf{x}}^{(t)}$ and predicts a distribution over unmasked completions:
\[
p_\theta\!\left(\mathbf{x} \mid \tilde{\mathbf{x}}^{(t)}\right)
= \prod_{i \in \mathcal{M}^{(t)}} p_\theta\!\left(x_i \mid \tilde{\mathbf{x}}^{(t)}\right),
\]
where $\mathcal{M}^{(t)}$ is the set of masked positions at step $t$.
The SFT objective is:
\[
\mathcal{L}_{\mathrm{SFT}}(\theta) =
-\mathbb{E}_{\mathbf{x}, t, \mathcal{M}}\!\left[
  \sum_{i\in\mathcal{M}} \log p_\theta(x_i \mid \tilde{\mathbf{x}}^{(t)})
\right].
\]
The RL stage optimises a reward $r(\mathbf{x})$ that combines generation quality
(from a judge model) and inference efficiency (number of denoising steps to convergence):
\[
\mathcal{L}_{\mathrm{RL}}(\theta) =
-\mathbb{E}_{\mathbf{x}\sim p_\theta}\!\left[r(\mathbf{x})\right]
+ \lambda\,D_{\mathrm{KL}}(p_\theta \| p_{\mathrm{ref}}),
\]
where $p_{\mathrm{ref}}$ is the SFT checkpoint and $\lambda$ controls KL regularisation.

\section{Learning or Inference Procedure}

\textbf{Stage 1---SFT.}
Starting from the Gemma~4 AR checkpoint, the model is fine-tuned on text data using
the masked denoising objective above. Less than 10\% of the original pre-training
token budget is required, as the AR checkpoint provides strong representations that
transfer directly to the bidirectional denoising task.

\textbf{Stage 2---RL + Sampler Distillation.}
A combined training step jointly (i) uses RL to improve response quality by rewarding
high-judge-score outputs and (ii) distils a fast sampler that fills the canvas in
fewer denoising steps without quality loss.

\textbf{Inference.}
The model is applied with block-autoregressive sampling: a canvas of 256 masked tokens
is initialised, then denoised iteratively. Generation can be quantised to 4-bit
precision to fit within 18\,GB VRAM.

\section{What the Guarantee Says}

The two-stage recipe inherits Gemma~4's general language understanding and instruction
following, while the RL stage provides an explicit optimisation target linking generation
quality to reward. The KL regularisation in the RL objective bounds the distance from the
SFT checkpoint, preventing degenerate solutions that score high on the judge but collapse
the output distribution. The sampler distillation guarantee is that the fast sampler matches
the full sampler's output distribution within a KL tolerance tuned during training.

\section{Experimental Findings}

\begin{itemize}[itemsep=2pt,parsep=0pt]
  \item \textbf{1000+ tokens/second} on a single NVIDIA H100 (batch size~1, int4 quantisation).
  \item \textbf{700+ tokens/second} on a consumer NVIDIA GeForce RTX~5090.
  \item \textbf{Up to $4\times$ faster} than the Gemma~4 AR baseline on dedicated GPU hardware.
  \item Fits in \textbf{18\,GB VRAM} when quantised to 4-bit, enabling local deployment
    on high-end consumer GPUs without multi-GPU setup.
  \item Quality competitive with the equivalent AR model on standard instruction-following
    and language-understanding benchmarks.
\end{itemize}

\section{Ablations and Interpretation}

\textbf{Stage~1 only vs.\ Stage~1+2.} The SFT-only checkpoint has higher variance
in denoising step count; Stage~2 distillation produces a sampler that converges in
fewer iterations, explaining the efficiency gains beyond what the base denoiser achieves.

\textbf{Canvas size.} The 256-token canvas size balances between parallelism (larger
is better) and the difficulty of predicting many tokens at once with high confidence
(larger is worse). Ablations showed 256 to be near-optimal for the 26B parameter scale.

\textbf{MoE vs.\ dense.} The 26B-A4B MoE architecture activates fewer parameters per
token than a dense 4B model would but leverages a much larger total capacity. The sparse
activation is essential for hitting the latency targets at this parameter scale.

\section*{Reference}

DiffusionGemma Team (Adrien Ali Ta\"iga et al.). \textbf{DiffusionGemma Technical Report}.
arXiv:2608.00146, August 2026.
\url{https://arxiv.org/abs/2608.00146}

\end{document}
